\section{Introduction}
\label{sec:introduction}

When you ask a language model how a refrigerator works, it tells you.
When you ask what the research says about whether gender is purely a social construct, it hedges.
Both are legitimate questions with substantive scientific answers, yet frontier LLMs treat them very differently --- not by refusing the second, but by wrapping it in qualifications, caveats, and appeals to complexity that dilute the response.
We call this phenomenon \emph{implicit avoidance}, and we show that it is widespread, consistent across model families, and distinct from the explicit refusal behaviors studied in prior work.

\para{Why this matters.}
LLMs are increasingly deployed in education, healthcare, legal advice, and scientific research --- domains where users seek clear, direct answers on difficult topics.
If models systematically hedge on questions involving social controversy, institutional critique, or epistemic uncertainty, they degrade information quality on precisely the topics where directness is most valuable.
Unlike explicit refusal, which is easy to detect and audit, implicit avoidance is subtle: the model \emph{does} answer, just poorly.

\para{The gap.}
A growing body of work studies LLM overrefusal --- the tendency to refuse safe prompts \citep{cui2024orbench, rottger2024xstest, xie2024sorrybench}.
These benchmarks measure binary outcomes (refuse vs.\ comply) and have revealed a near-universal tradeoff between safety and helpfulness \citep{cui2024orbench}.
However, they do not capture the \emph{implicit} avoidance we study here: hedging language, unnecessary disclaimers, and excessive qualification that reduce response quality without triggering a refusal.
Separately, work on emergent misalignment \citep{betley2025emergent} and persona vectors \citep{chen2025persona} has shown that finetuning can induce broad behavioral shifts through implicit persona generalization.
No prior work has connected these lines of inquiry to ask whether implicit avoidance arises from the same mechanism.

\para{Our approach.}
We design a behavioral probing experiment that measures implicit avoidance across a spectrum of topic types.
We construct 100 prompts in five categories: safe controls, known safety violations, and three types of \grayzones --- social taboos, institutional critique, and epistemic discomfort.
We query five frontier LLMs (\gptfour, \claude, \deepseek, \llama, \gemini) and score responses using both heuristic metrics (hedging density, disclaimer density) and an LLM judge (\textsc{GPT-4.1-mini}) on six avoidance dimensions.
We then analyze cross-model correlations to test whether avoidance patterns are shared across organizations.

\para{Key results.}
Gray zone topics elicit significantly higher implicit avoidance than safe controls (Cohen's $d = 1.35$, $p < 10^{-6}$), manifesting as $5.6\times$ more hedging and $3.5\times$ more disclaimers, with near-zero explicit refusal.
Per-prompt avoidance scores are highly correlated across the four models in our main analysis (mean Spearman $\rho = 0.78$), suggesting shared origins in training data rather than independent safety tuning decisions.
Epistemic discomfort topics --- where being wrong has social consequences --- trigger the strongest avoidance (mean score 1.35 vs.\ 0.91 for institutional critique and 0.82 for social taboos).

\para{Contributions.}
\begin{itemize}[leftmargin=*,itemsep=0pt,topsep=0pt]
    \item We identify and operationalize \emph{implicit avoidance} as a failure mode distinct from explicit refusal, measured through hedging, disclaimers, and deflection signals.
    \item We construct a 100-prompt behavioral probing suite with five topic categories that separates gray zone avoidance from both safe baselines and known safety refusals.
    \item We demonstrate that implicit avoidance is highly correlated across four frontier LLMs from different organizations ($\rho = 0.78$), providing evidence for persona generalization from shared training corpora.
    \item We show that epistemic discomfort --- not social taboo --- is the strongest trigger for implicit avoidance, suggesting that models are most cautious when being wrong carries social risk.
\end{itemize}
